\naslov{Rodovne funkcije}

\begin{definicija}
  \pojem{Rodovna funkcija} nenegativne celoštevilske slučajne spremenljivke $X$
  je dana z $G_X(s) = \sum_{k=0}^\infty s^k P(X=k)$.
\end{definicija}

Potenčna vrsta konvergira enakomerno na $[-1, 1]$, ker je tam dominirana z vrsto
$\sum_{k=0}^\infty P(X=k) = 1$.
Torej je $G_X(s)$ zvezna na $[-1, 1]$.
Na intervalu $(-1, 1)$ je vrsta neskončnokrat odvedljiva.
Iz formule $E(f(x)) = \sum_k f(k) P(X=k)$ za $f(x) = s^x$ dobimo
\[
  E(s^X) = \sum_{k=0}^\infty s^k P(X=k) = G_X(s).
\]

\begin{izrek}
  Naj bosta $X, Y$ neodvisni spremenljivki.
  Velja $G_{X+Y}(s) = G_X(s) G_Y(s)$.
\end{izrek}

\begin{proof}
  Sledi iz dejstva $E(s^{X+Y}) = E(s^X s^Y)$.
\end{proof}

\vprasanje{Povej nekaj lastnosti rodovnih funkcij.}

\begin{izrek}
  Naj bo $X$ slučajna spremenljivka. Potem velja
  \begin{itemize}
  \item $E(X) = \lim_{s \uparrow 1} G_X'(s)$
  \item $E(X (X-1) (X-2) \ldots (X-k+1)) = \lim_{s \uparrow 1} G_X^{(k)}(s)$
  \end{itemize}
\end{izrek}

\begin{proof}
  Samo za prvo točko.
  Opazimo, da ima $G_X$ pozitivne koeficiente, torej imajo vsi odvodi pozitivne
  koeficiente.
  Sledi, da so odvodi na $(0,1)$ povsod naraščajoči.

  Za začetek predpostavimo $E(X) < \infty$.
  Za $s \in (0,1)$ velja
  \[
	\sum_{k=1}^N k P(X = k) s^{k-1} \le G_X'(s) \le E(X),
  \]
  kjer prva neenakost velja, ker je vrsta na levi glava vrste za $G_X'$.
  Funkcija $G_X'$ je na $(0,1)$ naraščajoča in omejena, torej limita $\lim_{s
	\uparrow 1} G_X'(s)$ obstaja.
  Levi člen v zgornji neenakosti konvergira k $E(X)$ za $N \to \infty$, torej
  trditev velja po izreku o sendviču.

  Če je $E(X) = \infty$, potem levi člen v zgornji neenačbi divergira za $N \to
  \infty$, in je posledično tudi
  \[
	\lim_{s \uparrow 1} G_X'(s) = \infty.
  \]
\end{proof}